\chapter{Bibliographie et webographie}

\begin{thebibliography}{99}

\bibitem{express}
OpenJS Foundation.
\newblock \textit{Express — web framework for Node.js}.
\newblock \url{https://expressjs.com/}, consulté en 2026.

\bibitem{flutter}
Google LLC.
\newblock \textit{Flutter documentation}.
\newblock \url{https://docs.flutter.dev/}, consulté en 2026.

\bibitem{mysql}
Oracle Corporation.
\newblock \textit{MySQL 8.0 Reference Manual}.
\newblock \url{https://dev.mysql.com/doc/refman/8.0/en/}, consulté en 2026.

\bibitem{jwt}
IETF.
\newblock RFC 7519 — JSON Web Token (JWT).
\newblock \url{https://www.rfc-editor.org/rfc/rfc7519}, 2015.

\bibitem{oauth}
IETF.
\newblock RFC 6749 — The OAuth 2.0 Authorization Framework.
\newblock \url{https://www.rfc-editor.org/rfc/rfc6749}, 2012.

\bibitem{react}
Meta Platforms, Inc.
\newblock \textit{React — A JavaScript library for building user interfaces}.
\newblock \url{https://react.dev/}, consulté en 2026.

\bibitem{vite}
Vite Team.
\newblock \textit{Vite — Next Generation Frontend Tooling}.
\newblock \url{https://vitejs.dev/}, consulté en 2026.

\bibitem{uml}
Object Management Group (OMG).
\newblock \textit{Unified Modeling Language}.
\newblock \url{https://www.omg.org/spec/UML/}, consulté en 2026.

\end{thebibliography}

\section*{Webographie complémentaire}
\addcontentsline{toc}{section}{Webographie complémentaire}

\begin{itemize}[leftmargin=*]
  \item Documentation Node.js : \url{https://nodejs.org/docs/}
  \item Documentation Dart : \url{https://dart.dev/guides}
  \item Référence PlantUML : \url{https://plantuml.com/}
  \item GitHub Actions : \url{https://docs.github.com/actions}
\end{itemize}
